\chapter{Conclusions and Future Scope of Study}
\label{ch:conclusions}

\section{Summary of the Work}
This project built a simulation framework for federated learning over a UAV swarm
whose updates travel over lossy air-to-ground links and in which some UAVs are
Byzantine. The framework supports hierarchical aggregation at three base stations,
handover, packet erasures from a synthetic jammed channel or a measured UAV trace,
five attacks and twelve selectable aggregation rules (FedAvg, the CosAvg-style
baseline, median, trimmed mean, Krum, Multi-Krum, FLTrust, three CCRA variants and the
base paper's two erasure rules), with automated tests. The controlled study evaluated
six methods; a separate \BpNumRuns-run follow-up evaluated the two base-paper rules.
A FedAvg
baseline was established first. A reputation-based aggregator, CCRA, was then
designed with two forms of channel awareness, erasure masking and telemetry-weighted
reputation. After a promising but inconclusive single-seed study, a prespecified,
paired study of \NumRuns{} runs separated the two forms and measured each one's
effect. The results have been written up as a manuscript in IEEE letter format.

\section{Conclusions}
\begin{enumerate}
\item Non-IID data reduced FedAvg accuracy from 62.6\% to 60.2\% in the baseline,
  and a 30\% label-flip attack reduced it by 8.9 points on the ideal channel and
  7.8 points on the jammed channel in the preliminary study (last-five-round means).
\item None of the four aggregators that judge updates against a trusted reference
  update diverged in its 120 runs (480 runs combined), whereas FedAvg and the
  coordinate-wise median diverged in \DivFedavg{} and \DivMedian{} of 120 runs each.
\item Ignoring lost coordinates instead of treating them as zeros (erasure masking)
  improved accuracy under attack on a 5\,dB link by \MaskGain{} percentage points
  (95\% CI \ci{\MaskLo}{\MaskHi}).
\item Using link telemetry to slow reputation rewards had a meaningful negative
  effect: the prespecified primary estimate was $\PrimGain$ points
  (\ci{\PrimLo}{\PrimHi}), accompanied by a rise in honest-update rejection from
  \RejMask\% to \RejCcra\%.
\item A plausible cause is structural. With slowed rewards but full-rate penalties,
  the constant-evidence drift analysis~\eqref{eq:fixed-point} shows that weaker
  evidence lowers the equilibrium of an honest UAV's reputation from its mean score
  toward the lower tail of its scores, and the measured gate decisions are
  consistent with this mechanism.
\item On weak links, the exploratory comparison of the best variant with FLTrust was
  inconclusive; a clear advantage over FLTrust appeared only at 15\,dB.
\item Against the base paper, mask-only CCRA was \VsBpAttack{} points more accurate
  than its better erasure rule under attack at 5\,dB (\ci{\VsBpAttackLo}{\VsBpAttackHi}),
  with a signed clean contrast of $\VsBpClean$ points (mask-only minus fresh-only;
  \ci{\VsBpCleanLo}{\VsBpCleanHi}). FLTrust also gains \FltVsBpAttack{} points
  over fresh-only under attack. This large gain reflects trusted-reference defence;
  our distinct findings concern masking and the telemetry rule.
\item Within the base paper's rules, fresh-only exceeds reuse-old by
  \BpCleanGain{} points on clean stressed links
  (\ci{\BpCleanLo}{\BpCleanHi}). The rules diverged in \BpDivFresh/120 and
  \BpDivReuse/120 runs; neither survived the tested sign-flip attack. This does not
  contradict the original paper's honest-client convergence analysis.
\end{enumerate}

\section{Scope for Future Work (Phase II)}
\begin{enumerate}
\item \textbf{Extend the completed base-paper comparison.} Repeat key contrasts on
  a common hardware/software stack and fresh seeds. Evaluate whether filling lost
  packets from a trusted UAV's previous parameters can improve on masking without
  retaining poisoned values.
\item \textbf{Redesign the reward rule} so that weak links do not bias reputation
  downwards, for example by making the update symmetric in the evidence, or by
  normalizing each UAV's score against its own recent history, and freeze the rule
  before testing.
\item \textbf{Confirm on fresh seeds and conditions.} Evaluate the revised rule on
  new seeds and channel settings rather than the ten used here, so that tuning
  cannot inflate the result.
\item \textbf{Mobility and handover.} Use traces or mobility models with frequent
  handovers, and more base stations, to test whether telemetry helps when the
  dominant impairment is handover rather than fading.
\item \textbf{Stronger attackers.} Add adaptive attacks that target the reputation
  itself and attackers that spoof their link telemetry.
\item \textbf{Network-level simulation.} Replace the packet-error model with a
  network simulator such as NS-3 to capture medium access, retransmissions and
  latency.
\item \textbf{Generalization.} Repeat the key comparisons with another dataset, a
  larger model and a larger swarm, and submit the manuscript.
\end{enumerate}
